\documentclass{report}
\usepackage{luatexja} % LuaTeXで日本語を使うためのパッケージ
\usepackage{luatexja-fontspec} % LuaTeX用の日本語フォント設定

% --- 数学関連 ---
\usepackage{amsmath, amssymb, amsfonts, mathtools, bm, amsthm} % 基本的な数学パッケージ
\usepackage{type1cm, upgreek} % 数式フォントとギリシャ文字
\usepackage{physics, mhchem} % 物理や化学の記号や式の表記を簡単にする

% --- 表関連 ---
\usepackage{multirow, longtable, tabularx, array, colortbl, dcolumn, diagbox} % 表のレイアウトを柔軟にする
\usepackage{tablefootnote, truthtable} % 表中に注釈を追加、真理値表
\usepackage{tabularray} % 高度な表組みレイアウト

% --- グラフィック関連 ---
\usepackage{tikz, graphicx} % 図の描画と画像の挿入
\usepackage{background} % ウォーターマークの設定
\usepackage{caption, subcaption} % 図や表のキャプション設定
\usepackage{float, here} % 図や表の位置指定

% --- レイアウトとページ設定 ---
\usepackage{fancyhdr} % ページヘッダー、フッター、余白の設定
\usepackage[top = 20truemm, bottom = 20truemm, left = 20truemm, right = 20truemm]{geometry}
\usepackage{fancybox, ascmac} % ボックスのデザイン

% --- 色とスタイル ---
\usepackage{xcolor, color, colortbl, tcolorbox} % 色とカラーボックス
\usepackage{listings, jvlisting} % コードの色付けとフォーマット

% --- 参考文献関連 ---
\usepackage{biblatex, usebib} % 参考文献の管理と挿入
\usepackage{url, hyperref} % URLとリンクの設定

% --- その他の便利なパッケージ ---
\usepackage{footmisc} % 脚注のカスタマイズ
\usepackage{multicol} % 複数段組
\usepackage{comment} % コメントアウトの拡張
\usepackage{siunitx} % 単位の表記
\usepackage{docmute}
% \usepackage{appendix}
% --- tcolorboxとtikzの設定 ---
\tcbuselibrary{theorems, breakable} % 定理のボックスと改ページ設定
\usetikzlibrary{decorations.markings, arrows.meta, calc} % tikzの装飾や矢印の設定

% --- 定理スタイルと数式設定 ---
\theoremstyle{definition} % 定義スタイル
\numberwithin{equation}{section} % 式番号をサブセクション単位でリセット

% --- hyperrefの設定 ---
\hypersetup{
  setpagesize = false,
  bookmarks = true,
  bookmarksdepth = tocdepth,
  bookmarksnumbered = true,
  colorlinks = false,
  pdftitle = {}, % PDFタイトル
  pdfsubject = {}, % PDFサブジェクト
  pdfauthor = {}, % PDF作者
  pdfkeywords = {} % PDFキーワード
}

% --- siunitxの設定 ---
\sisetup{
  table-format = 1.5, % 小数点以下の桁数
  table-number-alignment = center, % 数値の中央揃え
}

% --- 透かし画像の設定 --- 
\backgroundsetup{
  scale=0.5,                       % 画像のスケール
  % color=black,                   % 画像の色（透かし用に半透明が推奨）
  opacity=0.2,                   % 透かしの透明度（0が完全透明、1が完全不透明）
  angle=0,                       % 画像の角度
  position = current page.south east,  % ページの右下
  hshift=-6cm, % 右方向へのシフト（負の値で内側に移動）
  vshift=5cm,  % 上方向へのシフト（正の値で内側に移動）
  contents={\includegraphics{./fig/appilogo-circular-full.png}} % 画像のパス
}

% --- その他の設定 ---
\allowdisplaybreaks % 数式の途中改ページ許可
\newcolumntype{t}{!{\vrule width 0.1pt}} % 新しいカラムタイプ
\newcolumntype{b}{!{\vrule width 1.5pt}} % 太いカラム
\UseTblrLibrary{amsmath, booktabs, counter, diagbox, functional, hook, html, nameref, siunitx, varwidth, zref} % tabularrayのライブラリ
\setlength{\columnseprule}{0.4pt} % カラム区切り線の太さ
\captionsetup[figure]{font = bf} % 図のキャプションの太字設定
\captionsetup[table]{font = bf} % 表のキャプションの太字設定
\captionsetup[lstlisting]{font = bf} % コードのキャプションの太字設定
\captionsetup[subfigure]{font = bf, labelformat = simple} % サブ図のキャプション設定
\setcounter{secnumdepth}{4} % セクションの深さ設定
\newcolumntype{d}{D{.}{.}{5}} % 数値のカラム
\newcolumntype{M}[1]{>{\centering\arraybackslash}m{#1}} % センター揃えのカラム
\DeclareMathOperator{\diag}{diag}
\everymath{\displaystyle} % 数式のスタイル
\newcommand{\inner}[2]{\left\langle #1, #2 \right\rangle}
\renewcommand{\figurename}{図}
\renewcommand{\i}{\mathrm{i}} % 複素数単位i
\renewcommand{\laplacian}{\grad^2} % ラプラシアンの記号
\renewcommand{\thesubfigure}{(\alph{subfigure})} % サブ図の番号形式
\newcommand{\m}[3]{\multicolumn{#1}{#2}{#3}} % マルチカラムのショートカット
\renewcommand{\r}[1]{\mathrm{#1}} % mathrmのショートカット
\newcommand{\e}{\mathrm{e}} % 自然対数の底e
\newcommand{\Ef}{E_{\mathrm{F}}} % フェルミエネルギー
\renewcommand{\c}{\si{\degreeCelsius}} % 摂氏記号
\renewcommand{\d}{\r{d}} % d記号
\renewcommand{\t}[1]{\texttt{#1}} % タイプライタフォント
\newcommand{\kb}{k_{\mathrm{B}}} % ボルツマン定数
% \renewcommand{\phi}{\varphi} % ϕをφに変更
\renewcommand{\epsilon}{\varepsilon}
\newcommand{\fullref}[1]{\textbf{\ref{#1} \nameref{#1}}}
\newcommand{\reff}[1]{\textbf{図\ref{#1}}} % 図参照のショートカット
\newcommand{\reft}[1]{\textbf{表\ref{#1}}} % 表参照のショートカット
\newcommand{\refe}[1]{\textbf{式\eqref{#1}}} % 式参照のショートカット
\newcommand{\refp}[1]{\textbf{コード\ref{#1}}} % コード参照のショートカット
\renewcommand{\lstlistingname}{コード} % コードリストの名前
\renewcommand{\theequation}{\thesection.\arabic{equation}} % 式番号の形式
\renewcommand{\footrulewidth}{0.4pt} % フッターの線
\newcommand{\mar}[1]{\textcircled{\scriptsize #1}} % 丸囲み文字
\newcommand{\combination}[2]{{}_{#1} \mathrm{C}_{#2}} % 組み合わせ
\newcommand{\thline}{\noalign{\hrule height 0.1pt}} % 細い横線
\newcommand{\bhline}{\noalign{\hrule height 1.5pt}} % 太い横線

% --- カスタム色定義 ---
\definecolor{burgundy}{rgb}{0.5, 0.0, 0.13} % バーガンディ色
\definecolor{charcoal}{rgb}{0.21, 0.27, 0.31} % チャコール色
\definecolor{forest}{rgb}{0.0, 0.35, 0} % 森の緑色

% --- カスタム定理環境の定義 ---
\newtcbtheorem[number within = chapter]{myexc}{練習問題}{
  fonttitle = \gtfamily\sffamily\bfseries\upshape,
  colframe = forest,
  colback = forest!2!white,
  rightrule = 1pt,
  leftrule = 1pt,
  bottomrule = 2pt,
  colbacktitle = forest,
  theorem style = standard,
  breakable,
  arc = 0pt,
}{exc-ref}
\newtcbtheorem[number within = chapter]{myprop}{命題}{
  fonttitle = \gtfamily\sffamily\bfseries\upshape,
  colframe = blue!50!black,
  colback = blue!50!black!2!white,
  rightrule = 1pt,
  leftrule = 1pt,
  bottomrule = 2pt,
  colbacktitle = blue!50!black,
  theorem style = standard,
  breakable,
  arc = 0pt
}{proposition-ref}
\newtcbtheorem[number within = chapter]{myrem}{注意}{
  fonttitle = \gtfamily\sffamily\bfseries\upshape,
  colframe = yellow!20!black,
  colback = yellow!50,
  rightrule = 1pt,
  leftrule = 1pt,
  bottomrule = 2pt,
  colbacktitle = yellow!20!black,
  theorem style = standard,
  breakable,
  arc = 0pt
}{remark-ref}
\newtcbtheorem[number within = chapter]{myex}{例題}{
  fonttitle = \gtfamily\sffamily\bfseries\upshape,
  colframe = black,
  colback = white,
  rightrule = 1pt,
  leftrule = 1pt,
  bottomrule = 2pt,
  colbacktitle = black,
  theorem style = standard,
  breakable,
  arc = 0pt
}{example-ref}
\newtcbtheorem[number within = chapter]{exc}{Requirement}{myexc}{exc-ref}
\newcommand{\rqref}[1]{{\bfseries\sffamily 練習問題 \ref{exc-ref:#1}}}
\newtcbtheorem[number within = chapter]{definition}{Definition}{mydef}{definition-ref}
\newcommand{\dfref}[1]{{\bfseries\sffamily 定義 \ref{definition-ref:#1}}}
\newtcbtheorem[number within = chapter]{prop}{命題}{myprop}{proposition-ref}
\newcommand{\prref}[1]{{\bfseries\sffamily 命題 \ref{proposition-ref:#1}}}
\newtcbtheorem[number within = chapter]{rem}{注意}{myrem}{remark-ref}
\newcommand{\rmref}[1]{{\bfseries\sffamily 注意 \ref{remark-ref:#1}}}
\newtcbtheorem[number within = chapter]{ex}{例題}{myex}{example-ref}
\newcommand{\exref}[1]{{\bfseries\sffamily 例題 \ref{example-ref:#1}}}
% --- 再定義コマンド ---
% \mathtoolsset{showonlyrefs=true} % 必要な式番号のみ表示
\pagestyle{fancy} % ヘッダー・フッターのスタイル設定
\chead{応用量子物性講義ノート} % 中央ヘッダー
% \rhead{}
\fancyhead[R]{\rightmark}
\renewcommand{\sectionmark}[1]{\markright{\thesection\ #1}}
\cfoot{\thepage} % 中央フッターにページ番号
\lhead{}
\rfoot{Yuto Masuda and Haruki Aoki} % 右フッターに名前
\setcounter{tocdepth}{4} % 目次の深さ
\makeatletter
\@addtoreset{equation}{section} % サブセクションごとに式番号をリセット
\makeatother

% --- メタ情報 ---
\title{応用量子物性講義ノート}
\date{更新日\today}
\author{Yuto Masuda and Haruki Aoki}

\begin{document}
  本節では，摂動の大きさが時間に依存しない場合を考える．
  摂動の大きさが時間に依存しないことは，定常摂動とは異なる．
  定常摂動では，量子系に摂動をかけ始めてから十分長い時間が経った状況のみ扱えるが，非定常摂動を考えると，量子系の時間変化を任意の時刻で追跡できる．
  \par
  さて，時刻$t = 0$から一定の摂動$\hat{V}$を加え始めたとき考える．$\hat{V}(t)$は，
  \begin{align}
    \hat{V}(t) =
    \begin{dcases}
      0 & t \leq 0 \\
      \hat{V} & t \geq 0
    \end{dcases}
  \end{align}
  と書ける．
  $V_{fi}(t)$は，
  \begin{align}
    V_{fi}(t) = \mel**{f}{\hat{V}(t)}{i} =
    \begin{dcases}
      0 & t \leq 0 \\
      \mel**{f}{\hat{V}}{i} & t \geq 0
    \end{dcases}
  \end{align}
  である．$V_{fi}$を
  \begin{align}
    V_{fi} \coloneqq \mel**{f}{\hat{V}}{i}
  \end{align}
  定義する．また，$t = 0$で量子状態は$\ket{i}$であったとする．
  時刻$t$で状態$\ket{f}$である確率を考える．
  \refe{interaction-picture-time-development-of-qc}の第2式より，
  \begin{align}
    c_{f, i}^{(1)}(t) &= -\frac{\i}{\hbar}\int_{0}^{t}V_{fi}\e^{\i\omega_{fi}t'}\dd{t'} \\
    &= -\frac{V_{fi}}{\hbar\omega_{fi}}\qty(\e^{\i\omega_{fi}t} - 1) \\ 
    &= -\frac{V_{fi}}{\hbar\omega_{fi}}\exp\qty(\frac{\i\omega_{fi}}{2}t)\qty[\exp\qty(\frac{\i\omega_{fi}}{2}t) - \exp\qty(-\frac{\i\omega_{fi}}{2}t)] \\ 
    &= 2\i\frac{V_{fi}}{\hbar\omega_{fi}}\exp\qty(\frac{\i\omega_{fi}}{2}t)\cdot\sin\qty(\frac{\omega_{fi}}{2}t)
  \end{align}
  と計算できるから，$\ket{i}$から$\ket{f}$への遷移確率$\abs{c_{f, i}^{(1)}(t)}^2$は，
  \begin{align}
    \abs{c_{f, i}^{(1)}(t)}^2 &= \frac{\abs{V_{fi}}^2}{\hbar^2}\qty[\frac{\sin\qty(\frac{\omega_{fi}}{2}t)}{\qty(\frac{\omega_{fi}}{2})}]^2 \\ 
    &= \frac{\abs{V_{fi}}^2}{\hbar^2}t^2\qty[\frac{\sin\qty(\frac{\omega_{fi}}{2}t)}{\qty(\frac{\omega_{fi}}{2}t)}]^2 \label{v-const-cfi-t}
  \end{align}
  である．\refe{v-const-cfi-t}において，$t$を固定して$\omega_{fi}$が満たすべき物理条件を考える．
  まず，$\abs{c_{f, i}^{(1)}(t)}^2$が0でない，つまり，$\ket{i}$から$\ket{f}$への遷移が起きるのは，$\sin\qty(\frac{\omega_{fi}}{2}t)$が0でないときである．
  また，$\qty[\frac{\sin\qty(\frac{\omega_{fi}}{2}t)}{\qty(\frac{\omega_{fi}}{2})}]^2$を$\omega_{fi}$の関数とみたとき，
  $\sin$が1周期の間のみで有効な値を持つ．
  言い換えれば，$\abs{\omega_{fi}} < \frac{2\pi}{t}$で有効な値をもつ．
  これは，$\sin$の2周期目の最大値が1周期目の最大値の$\qty(\frac{2}{3\pi})^2\simeq 4.4\ \%$程度であること依る．
  従って，$t$が小さいときはこの範囲が十分広く，$\omega_{fi} \neq 0$の状態への遷移が起こるが，
  $t$が大きいときは$\omega_{fi}\simeq 0$の状態($\sin$の1周期目)への遷移しか起きない．
  摂動を与えた瞬間に測定を行うことは困難であるという実験上の理由から，通常は$t \to \infty$としておけば良い．
  また，規格化されたデルタ関数$\tilde{\delta}$を用いると，$\r{sinc}$関数の性質から，
  \begin{align}
    \lim_{t\to\infty}\qty(\frac{\sin\qty(\frac{\omega_{fi} t}{2})}{\frac{\omega_{fi} t}{2}})^2 = \frac{2\pi}{t}\tilde{\delta}(\omega_{fi})
  \end{align}
  を使うと\footnote{$\delta(ax)=\frac{1}{\abs{a}}\delta(x)$}，
  \begin{align}
    \lim_{t\to\infty}\abs{c_{f, i}^{(1)}(t)}^2 &= \frac{\abs{V_{fi}}^2}{\hbar^2}2\pi t\delta(\omega_{fi})\\
    &= \frac{2\pi}{\hbar}\abs{V_{fi}}^2\delta(E_f - E_i)t
  \end{align}
  が得られる．$t$をあまりにも大きくすると，$\abs{c_{f, i}^{(1)}(t)}^2 > 1$となり，$\abs{c_{f, i}^{(1)}(t)}^2$の時刻$t$での$\ket{i}$から$\ket{f}$への遷移確率という
  物理的な意味が失われることに注意しながら，$t\to \infty$の領域での単位時間当たりの$\ket{i}$から$\ket{f}$への遷移確率$\omega_{i\to f}$は，
  \begin{align}
    \omega_{i\to f} &= \frac{\abs{c_{f, i}^{(1)}(t)}^2}{t} \\
    &= \frac{2\pi}{\hbar}\abs{\mel**{f}{\hat{V}}{i}}^2\delta(E_f - E_i)
  \end{align}
  である．これを\textbf{Fermiの黄金律}という\footnote{Enrico Fermi(1901-1954)}．
  \begin{itembox}[l]{Fermiの黄金律}
    \begin{align}
      \omega_{i\to f} = \frac{2\pi}{\hbar}\abs{\mel**{f}{\hat{V}}{i}}^2\delta(E_f - E_i)\label{fermis-golden-rule}
    \end{align}
  \end{itembox}
  \begin{myex}{電子の弾性散乱}{}
    弾性散乱では散乱前後で粒子のエネルギーが保存される．電子のエネルギーは
    \begin{align}
      E_{\bm{k'}} = \frac{\hbar^2\abs{\bm{k'}}^2}{2m}
    \end{align}
    と表される．よって，終状態は多く存在する．始状態$i$から終状態集合 $\qty{f}$への遷移確率を求める．
    \begin{align}
      \omega_{i \to \qty{f}} &= \sum_{f}\omega_{i \to f}\\
      &= \sum_{f}\frac{2\pi}{\hbar}\abs{\mel**{f}{\hat{V}}{i}}^2\delta(E_f - E_i)\\
      &= \frac{2\pi}{\hbar}\overline{\abs{\mel**{f}{\hat{V}}{i}}^2}\sum_{f}\delta(E_f - E_i)\\
      &= \frac{2\pi}{\hbar}\overline{\abs{\mel**{f}{\hat{V}}{i}}^2}\rho(E_f)
    \end{align}
    $\overline{\abs{\mel**{f}{\hat{V}}{i}}^2}$は散乱体の性質を，状態密度$\rho(E_f)$は物質の性質を反映している．
    この関係もFermiの黄金律という．
  \end{myex}
\end{document}
